\documentclass{article}
\usepackage{fortify}
\usepackage[active,tightpage]{preview}
\setlength\PreviewBorder{6pt}
\begin{document}
\begin{preview}
\begin{FortressCode}
\KWD{component} \TYP{NotationDualSumGenerator} \\
\KWD{import} \TYP{List}.\lbrace\ldots\rbrace \\
\KWD{export} \TYP{Executable} \\
 \\
\KWD{object} D(v\COLON\mathbb{R}64, d\COLON\mathbb{R}64) \\
\KWD{end} \\
\KWD{opr}\, \mathord{+}(x\COLONOP{}D, y\COLONOP{}D)\COLONOP{}D = D(x.v+y.v, x.d+y.d) \\
 \\
\KWD{object} \TYP{DSumReduction} \KWD{extends} \TYP{CommutativeMonoidReduction}\llbracket{}D\rrbracket \\
\2\+\VAR{empty}(\ultrathin)\COLONOP{}D = D(0.0,0.0) \\
  \VAR{join}(x\COLONOP{}D,y\COLONOP{}D)\COLONOP{}D = x+y\- \\
\KWD{end} \\
 \\
\FortressBeginFlushRightComment A one-argument overload avoids the zero-argument collision used by
generator-clause desugaring.  Test whether SUM <{$|$} ... {$|$}> can retain the
ordinary summation glyph for graph scalars. \FortressEndComment \\
\KWD{opr}\, \mathord{\sum}\bigl(g\COLON\TYP{Generator}\llbracket{}D\rrbracket\bigr)\COLONOP{}D = g.\VAR{reduce}(\TYP{DSumReduction}) \\
 \\
\VAR{run}(\ultrathin)\COLON(\ultrathin) =\; \KWD{do} \\
\2\+s\COLONOP{}D = \sum \bigl\langle\llbracket{}D\rrbracket\, D(1.0\, i, 1.0) \bigm| i \leftarrow 1\mathbin{\#}4\,\bigr\rangle \\
  \VAR{println}(s.v) \\
  \VAR{println}(s.d)\- \\
\KWD{end} \\
\KWD{end}
\end{FortressCode}
\end{preview}
\end{document}
